\documentclass[a4paper,12pt,numbers=enddot]{scrreprt}	% numbers=enddot fa seguire il puntino dopo il numero di cap. noenddot lo toglie 
\usepackage[utf8]{inputenc}	%imposta la codifica di input; invece di "latin1", anche "utf8" va bene
\usepackage[small]{caption}	%per i cara piccoli
\usepackage[english]{babel}	%l�ultima lingua � quella predefinita
\usepackage{listings}		%codice sorgente
\lstset{frame=shadowbox, rulesepcolor=\color{blue}, breaklines=true, basicstyle=\small\ttfamily, columns=fullflexible,  keepspaces=true}	%disegna un rettangolo attorno al codice sorgente
\usepackage{indentfirst}	%rientra il primo paragrafo di ogni sezione
\usepackage{amsmath}%,amssymb,amsthm       % indispensabili per la matematica
\usepackage{graphicx}		% per inserire le immagini
\usepackage{graphics}		%?
\usepackage{sidecap}		%per mettere caption a lato della figuta
\usepackage{caption}		% didascalie personalizzate
\usepackage[numbers]{natbib} %?
\usepackage{psfrag} %?
\usepackage{multirow}	%?
\usepackage{booktabs} 	%?
\usepackage{lscape}%?
\usepackage{longtable}%?
\usepackage{enumitem}%?
\usepackage{epsfig}		%per mettere fig usando comando tabular
\usepackage{makeidx}	%indice
\usepackage{varioref}
\usepackage[table]{xcolor}
\usepackage{epstopdf}
%Setta dimensioni body
\usepackage[total={16.5cm,25cm},
						top=3.5cm,
						bottom=3.5cm,
						left=2.5cm,
						right=2.5cm,
						%includefoot,
						centering]{geometry}

\usepackage{wallpaper}%?
\usepackage{color}%?
\usepackage{setspace}		%per definire con dopo l'interlinea con \onehalfspacing
\onehalfspacing
\usepackage{appendix}
\usepackage{acronym}%?
\usepackage{makecell}

%Intestazioni e pie� pag.  
\usepackage[automark,headsepline,footsepline,autooneside,plainheadsepline,plainfootsepline]{scrpage2}	%i comandi plain-ecccc servono a dire che nella prima pag. del nuovo cap. voglio queste cose.
\usepackage[pagebackref=true, colorlinks=true, linkcolor=blue, hyperindex=true, urlcolor=blue, citecolor = green, linktocpage=true]{hyperref}		% indici e riferimenti cliccabili
\pagestyle{scrheadings}
\clearscrheadfoot
\ihead [Stefano Bennati]{Stefano Bennati}	%Kopfzeile Innen
\chead {}					%Kopfzeile Mitte
\ohead [\headmark]{\headmark}			%Kopfzeile Au�en
\ifoot {}					%Fu�zeile Innen
\cfoot[\pagemark]{\pagemark} 			%Fu�zeile Mitte
\ofoot {}

\automark[chapter]{chapter}
\KOMAoptions{cleardoublepage=scrheadings}	%per lasciate style=scrheadings anche per le pag vuoto (svuotate col comando cleardoublepage), si puo� anche scegliere cleardoublepage=plain o empty

\footskip=30pt	 %	abbassa il margine inf della pag
\setlength{\headheight}{2\baselineskip}

\newcommand\SectionFontStyle{\rmfamily}			% geh�rt zu den beiden folgenden setkomafont , keine ahnung was es macht =)
\setkomafont{chapter}{\huge\bfseries\SectionFontStyle}	% Chapter in der gleichen schriftart SectionFontStyle
\setkomafont{sectioning}{\bfseries\SectionFontStyle}
%Crea i link nel pdf (vedere come si fa a farli verdino e con i collegamenti solo sui numeri, non questi rettangoli cessi rossi	

\makeindex
\begin{document}

\renewcommand{\bibname}{References}	%to change the name: bibliography to references

%\pagenumbering{roman}
%\pagestyle{plain}
%******************************************************************
% Materiale iniziale
%******************************************************************

\begin{titlepage}   %Frontespizio
\begin{center}
{\Large UNIVERSIT\`A DEGLI STUDI DI BRESCIA}

\vspace{0.2cm}

{\Large FACOLT\`A DI INGEGNERIA}

\begin{figure*}[htbp] 
\begin{center}
 \begin{tabular}{c c}
\raisebox{-.5\height}{\includegraphics[width=0.20\textwidth]{images/00-loghi/logo_UniBS.eps}}
 & \raisebox{-0.57\height}{\includegraphics[width=0.20\textwidth]{images/00-loghi/freiburg.png}}
\end{tabular}
\end{center}
\end{figure*}

% \vspace{0.5cm}

CORSO DI LAUREA IN INGEGNERIA INFORMATICA \\

TESI DI LAUREA SPECIALISTICA \\

%\vspace{0.8cm}
\vspace{1.5cm}
{\normalsize \emph{\textbf{ANALISI DI FATTIBILIT\`A ED IMPLEMENTAZIONE}}} \\

\vspace{0.2cm}

{\normalsize \emph{\textbf{DI UN SISTEMA DI ROBOTICA COGNITIVA}}} \\

\vspace{0.2cm}

{\normalsize \emph{\textbf{PER COMPITI DI NAVIGAZIONE}}} \\

%{\normalsize \emph{\textbf{IN REGIME INTERMITTENTE}}} \\

\vspace{0.5cm}

{\small \emph{\textbf{(Analysis of preconditions and implementation of a \\ Cognitive Robotic System for navigation tasks)}}} \\

\vspace{1.5cm}

\begin{flushleft}

{\normalsize Relatore:}
\begin{quote}
{\large \textbf{Ch.mo Prof. Riccardo Cassinis}} \\
\vspace{0.2cm}
\end{quote}

{\normalsize Correlatore:}
\begin{quote}
{\large \textbf{Ch.mo Prof. Marco Ragni}} \\
\vspace{0.2cm}
\end{quote}

\end{flushleft}

% \vspace{0.5cm}

\begin{flushright}
Laureando: \\
\textbf{Stefano Bennati} \\
\vspace{0.05cm}
\textbf{Matricola 72363} \\
\vspace{1.5cm}
\end{flushright}

\large{Anno Accademico 2010-2011}

\end{center}
\end{titlepage}

\pagenumbering{arabic}
\tableofcontents
\listoffigures

\newpage
\thispagestyle{empty}
\mbox{}

%\documentclass[a4paper,12pt]{book}
\begin{document}
\chapter{Introduction}
\section{Background}
\section{Problem description}
\section{Working Instruments}
\subsection{ACT-R}
\subsection{OpenCV}
	\input{./01/opencv/opencv}
\end{document}
\begin{tiny}
•
\end{tiny}

\pagenumbering{roman}

%\input{Chapters/1}
\newpage
\clearpage
\appendix
%\input{Chapters/appendice}
\newpage
\clearpage



%******************************************************************
% Materiale Finale
%******************************************************************
%******************************************************************
%References (two methods in Latex, this (bibTex) with the databe is the best
%******************************************************************
\cleardoublepage
\phantomsection	%gives hyperref something to hold on to when making the link.
\addcontentsline{toc}{chapter}{References}
\bibliographystyle{alpha}				%scelgo come formattare la biblio (abbrv abbrevia i nomi propri. plain non li abbrevia. Ce ne sono un sacco).
\bibliography{BeginningEnding/database_references.bib} 	%database dove prendere le info  
\nocite{*}						%stampa TUTTO quello presente nella banca dati, altrimenti solo la roba citata nella tesi

%NB in kile bisogna aggiungere BibTex al QuickBuild per far compilare la bibliografia. Metterlo subito dopo Latex.


\newpage										
\clearpage


\end{document}
